\section{Introduction}

The digitization of documents has become increasingly prevalent in various domains, including libraries, archives and businesses. The goal is to revitalize the information contained in the physical documents and make it more accessible and searchable in digital form as well as preserving it for future generations. This process involves scanning physical documents in order to create a digital representation of them. However, the quality of the scanned images can vary significantly due to factors such as lighting conditions, document physical state and the scanning equipment used. As a result, the scanned images may contain noise, shadows, and other artifacts that can hinder the process of the recognition using OCR (Optical Character Recognition) systems~\cite{Souibgui_2022}.

To address such challenges, preprocessing techniques are often applied beforehand in order to enhance the quality of the scanned images and improve the accuracy of OCR systems~\cite{GUPTA2007389}. One of the most common preprocessing techniques is image binarization, which converts a colorful or a greyscale image into a binary image consisting solely of black and white pixels~\cite{10.1007/s42979-020-00176-1}. According to Gupta et al., this process allows OCR systems to focus on the relevant information in the document and ignore the background, making the OCR systems more accurate~\cite{GUPTA2007389}.

While there are many algorithms for document or image binarization available, such as Niblack's approach and Sauvola's method, these algorithms can be computationally intensive as they often require processing each pixel or each cluster of pixels individually, which can be time-consuming, especially when working with large documents with high DPI~\cite{10.5120/3877-5418,10.1007/s11042-014-2269-7}. Therefore, there is a growing need for high-performance solutions that entail the usage of parallel processing techniques either on the CPU level or the GPU level. Furthermore, implementing parallelization on the CPU level can be achieved by splitting the work among many worker threads, whereas with GPU-level parallelization, the work can be distributed among many GPU threads, which are usually more efficient than CPU threads when it comes generally to image processing and rendering tasks~\cite{ansari2025acceleratingmatrixmultiplicationperformance}. This statement was shown by Ansari et al.\ as they did a performance comparison between parallelized CPU and GPU performance for large metrics and the GPU performance was consistently faster than the CPU by many folds. The reason for that is because GPUs tend to have significantly more cores than CPUs, allowing them to handle a larger amount of smaller computations in parallel.